\section{Introduction}
\label{sec:introduction}

Fix one viscosity \(\nu>0\), one divergence-free datum
\(u_0\in H^\infty(\R^3;\R^3)\), and one prescribed source
\[
 f\in C^\infty([0,\infty);\mathcal S(\R^3;\R^3)).
\]
Forced local theory produces a unique pressure-normalized classical history
\((u,p)\) on a maximal interval \([0,T_*)\). This paper follows that one
history. Its question is whether a finite value of \(T_*\) can terminate the
solution generated by \((u_0,\nu,f)\).

\begin{maintheorem}
\label{thm:main-global}
For every \(\nu>0\), every \(u_0\in\HsSigma(\R^3)\), and every
\(f\in C^\infty([0,\infty);\mathcal S(\R^3;\R^3))\), there is a unique
pressure-normalized pair
\[
 (u,p)\in C^\infty([0,\infty)\times\R^3;\R^3\times\R)
\]
satisfying
\[
 \partial_tu+(u\cdot\nabla)u+\nabla p=\nu\Delta u+f,
 \qquad \nabla\cdot u=0,
 \qquad u(0)=u_0,
 \qquad \lim_{\abs{x}\to\infty}p(t,x)=0.
\]
It obeys the forced energy identity
\[
 \frac12\norm{u(t)}_{L^2}^2
 +\nu\int_0^t\norm{\nabla u(s)}_{L^2}^2\dd s
 =\frac12\norm{u_0}_{L^2}^2
 +\int_0^t\ip{f(s)}{u(s)}_{L^2}\dd s
 \qquad(t\ge0).
\]
\end{maintheorem}

\begin{corollary}[Weak--strong uniqueness and pressure normalization]
\label{cor:main-weak-strong}
For the solution of the preceding theorem, every forced Leray--Hopf pair
\((v,q)\) with the same datum, viscosity, and
prescribed source has \(v=u\). Distributionally,
\(q=p+c(t)\) for a scalar distribution \(c\); if
\(q(t,\cdot)\in C_0(\R^3)\) almost everywhere, then \(q=p\) almost
everywhere.
\end{corollary}

Leray's weak solutions establish global finite-energy existence
\cite{Leray1934}. The regularity results of Serrin
\cite{Serrin1962}, Caffarelli--Kohn--Nirenberg \cite{CKN1982}, and
Escauriaza--Seregin--\v{S}ver\'ak \cite{ESS2003} identify increasingly sharp
ways in which a singularity would have to evade classical control.
Fujita--Kato and Kato provide the local strong theory used both at the
initial time and at the terminal restart \cite{FujitaKato1964,Kato1984}.
The proof below works directly with the complete derivative structure of the
prescribed forced history.

Forcing changes every level of that structure. Its solenoidal part
\(g=\Pp f\) enters the projected velocity equation. Its gradient part is
absorbed by the normalized pressure,
\[
 -\Delta p=\partial_i\partial_j(u_i u_j)-\nabla\cdot f,
 \qquad
 p=R_iR_j(u_i u_j)-(-\Delta)^{-1}\nabla\cdot f.
\]
Every temporal row contains the corresponding derivative
\(F_n=\partial_t^nf\); every energy row contains force work; and a restart at
\(T_*\) retains the same absolute-time source. These are parts of a single
forced velocity--pressure--source history.

The proof has two constructions from the same pressure-complete mixed jet.
For every finite pair of temporal and spatial orders, the datum-generated
branch produces a compatible whole-terminal norm rectangle containing the
selected rows and their adjacent rows.  That rectangle is the quantitative
mechanism: it keeps every selected derivative norm finite through the chosen
finite endpoint.  Restriction compatibility assembles the cofinal family of
finite rectangles into the complete space--time intersection.  Independently,
the critical-height tower and the positive-sign viscous diagonal relation
supply the relationship identities on the same coordinates.  Pairing,
squaring, and integrating the complete forced rows keeps every face of each
rectangle inside one finite estimate.  Because every mixed derivative belongs
to a finite rectangle, the adjacent rows have one common smooth endpoint.
The pressure formula and forced recurrences determine
the complete endpoint jet, and local theory restarts the equation from
\(u_*=u(T_*)\) with the same source. Overlap uniqueness then contradicts a
finite maximal time.

Section~\ref{sec:forced-argument-overview} derives this architecture in one
view. The analytic framework and the mixed-jet construction follow, then the
independent whole-terminal intersection and rectangle branches, their
composition, and the endpoint
continuation. The remaining
sections record matrix formulations and quantitative consequences of the same
finite rectangles.
